\section{Completion sets of minimum partial homomorphisms}\label{norm391:section}
In this section, we prove the bounds~\eqref{ov:normalized} on optimal covers (Theorem~\ref{norm391:normalized-cover}), which we use in Sections~\ref{sec:spantwo}--\ref{sec:tworeturns}. By Lemma~\ref{general:max-deletion}, this reduces to the following problem. We are given a measure on minimum partial homomorphisms, and we wish to find a minimum partial homomorphism whose completion sets have small total multiplicity. We solve this problem, for the orientations of interest, in Theorem~\ref{norm391:endpoint}.

The proof goes as follows. Suppose that every minimum partial homomorphism has completion sets of large total multiplicity. We first rescale the weights so that every partial homomorphism in the support of the measure meets the set $Z$ of vertices of weight zero (Lemma~\ref{norm391:cutoff}). We then modify minimum partial homomorphisms by putting a vertex at the deleted position and deleting one of the endpoints instead, which moves the deleted position by one. The new partial homomorphism must again have completion sets of large multiplicity, and this restricts how the positions mapped into $Z$ can be arranged. We exclude, in turn, partial homomorphisms with exactly one, two and three positions in $Z$, and then those with more, and so reach a contradiction.

\medskip

Throughout this section, $H$ is an orientation of a cycle of even order $r\ge4$, $D$ is a digraph, and the vertices of $D$ have nonnegative weights, which we call \emph{costs} and denote by $c$. We let $Z$ be the set of vertices of cost zero.
For brevity, we call partial homomorphisms into $D$ \emph{maps}, and we call homomorphisms of $H$ into $D$ \emph{full maps}. A \emph{minimum map} is a minimum partial homomorphism, that is, a map whose cost is minimum among all maps, over all deleted positions; this includes maps outside the support of any measure we consider.
For a map $Q$, we write $X_Q,Y_Q$ for
its completion sets. We consider a measure $\rho$
of mass $m>0$ on maps, with multiplicity function $p$ and total multiplicity
$P=(r-1)m$, and we write $\Phi_p(Q)=p(X_Q)+p(Y_Q)$ as before. We also put
\[
\nu=r/2-2.
\]
A position of a map is a \emph{zero position} if it is mapped into $Z$, and an endpoint of a map is a \emph{zero endpoint} if it is a zero position.

We want to show that some minimum map $Q$ has $\Phi_p(Q)\le(r-2)m=P-m$. As there are finitely many maps, if this fails, then there is some $\varepsilon>0$ such that
\begin{equation}\label{norm391:counter}
 \Phi_p(Q)\ge P-m+\varepsilon\quad\text{for every minimum map }Q.
\end{equation}
Most of this section consists of showing that~\eqref{norm391:counter} is impossible.
Suppose that~\eqref{norm391:counter} holds and that the completion sets of every minimum map are disjoint. We say that a set $A$ is \emph{avoided} if the completion pair of some minimum map avoids $A$. We use two simple consequences of~\eqref{norm391:counter} repeatedly. First, if $A$ is avoided, then $p(A)\le m-\varepsilon$. Indeed, the completion sets of the corresponding minimum map are disjoint and miss $A$, so their total multiplicity is at most $P-p(A)$. Second, the completion pair of every minimum map has multiplicity at least $p(Z)-m+\varepsilon$ on $Z$, as at most $P-p(Z)$ of its multiplicity lies outside $Z$.
Here $p$ is always the multiplicity function of $\rho$, also when we apply these bounds to minimum maps outside its support.

Two simple operations on a map are used throughout. Let $Q$ be a map with deleted position $j$. \emph{Inserting} a vertex $v$ of a completion set of $Q$ means putting $v$ at position $j$; the edge between $j$ and the corresponding endpoint is then satisfied. If we also delete the other endpoint, we obtain a new map, whose cost is the cost of $Q$ plus $c(v)$ minus the cost of the deleted endpoint. In particular, if $Q$ is minimum, $v\in Z$, and the deleted endpoint has positive cost, then the new map is cheaper than $Q$, which is impossible. Here and below, the \emph{completion set at an endpoint} is the completion set defined by the edge between this endpoint and the deleted position. We record the consequence.

\begin{observation}\label{insertion:rule}
If a minimum map has a zero endpoint and an endpoint of positive cost, then the completion set at its zero endpoint is disjoint from~$Z$.
\end{observation}

For the orientations of interest to us, the following conditions on $H$ hold, where $\Delta$ is the maximum drop of the word of $H$ (see Section~\ref{compact390:structural-section}):
\begin{equation}\label{norm391:range}
 r\ \text{is even},\qquad \text{the sign sum of }H\text{ is }2,\qquad
 \Delta\in\{2,3\},\qquad r\ge\Delta+6.
\end{equation}
We assume~\eqref{norm391:range} only where we say so.

\subsection{Height maps}
If a position $j$ is deleted from $H$, then the height walk of the path $H-j$ is obtained by reading the word of $H$ from $j+1$ to $j-1$, and we call its span the \emph{span of $H-j$}.

\begin{lemma}
\label{norm391:height}
Let $H$ be a cycle orientation with sign sum two and maximum drop $\Delta\ge1$. Then the following hold.
\begin{props}
\item\label{norm391:height:1} For every position $j$, the span of $H-j$ is at least $\Delta$, and, for some position $j$, it is at most $\Delta+1$.
\item\label{norm391:height:2} Let $z_0\to z_1\to\cdots\to z_{\Delta+1}$ be a directed walk, and let $1\le t\le\Delta$. Then there is an internal position $j$ of $H$, whose two incident edges are forward edges, such that $H-j$ has a homomorphism to this walk which maps both endpoints to $z_t$. In particular, the completion sets of this map are $N^-(z_t)$ and $N^+(z_t)$.
\item\label{norm391:height:3} Every tournament on $\Delta+3$ vertices receives a homomorphism from $H$.
\end{props}
\end{lemma}
\begin{proof}
Fix a proper cyclic interval of sign sum $-\Delta$; its complement has sign sum $\Delta+2$. Deleting a position inside the complement leaves the interval intact, and deleting a position inside the interval leaves the complement intact. Deleting one of the two common endpoints of the interval and its complement removes at most one edge from the complement. In each of these three cases, the span of $H-j$ is at least $\Delta$.

For the upper bound in~\ref{norm391:height:1}, rotate the word so that it starts at a position of minimum height, which we take to be zero, so that it ends at height two. Let $j$ be the last position at height zero. The two steps after $j$ are positive, as otherwise the walk would return to height zero. As the drop is at most $\Delta$, the heights before $j$ lie in $\{0,\ldots,\Delta\}$ and the heights after $j$ lie in $\{1,\ldots,\Delta+2\}$. Delete $j+1$, and read the remaining path from $j+2$ to the end, and then from the start to $j$, translating the latter heights by two. All heights of this path lie in $\{1,\ldots,\Delta+2\}$, and both of its endpoints have height two. Thus its span is at most $\Delta+1$.

For~\ref{norm391:height:2}, rotate the word instead so that it starts at the last position of minimum height, so that its heights are $s_0=0$, $s_r=2$ and $s_i>0$ for $0<i<r$. Let $K=\max_is_i$. An interval from a highest position to the end has drop $K-2$, so $K\le\Delta+2$. An interval which does not wrap around has drop at most $K-1$, and one which wraps around has drop at most $K-2$, so $K\ge\Delta+1$. In particular, the walk visits height $t+1$, and immediately before its first visit to $t+1$ it visits $t-1$ and $t$. Delete this position at height $t$, whose two incident steps are positive. Reading the remaining path from its successor to the end, and then from the start translated by two, gives heights in $\{1,\ldots,\Delta+2\}$ with both endpoints at height $t+1$. Subtract one, and send each position at height $i$ to $z_i$. Every edge of $H$ goes from a position at some height to a position at the next height, so this is a homomorphism, and both endpoints are mapped to $z_t$.

For~\ref{norm391:height:3}, note that the chosen interval of sign sum $-\Delta$ stays between its initial height and $\Delta$ below it. Similarly, its complement stays between $0$ and $\Delta+2$ above its own initial height, as an excursion beyond these ranges would give a larger drop. Therefore, sending heights to vertices of two directed paths of lengths $\Delta+2$ and $\Delta$ with common ends gives a homomorphism from $H$ to $D_{\Delta+2,\Delta}$. By Lemma~\ref{general:tournament-paths}, every tournament on $\Delta+3$ vertices receives a homomorphism from $D_{\Delta+2,\Delta}$, and composing these maps proves~\ref{norm391:height:3}.
\end{proof}

\subsection{Restricting to maps through zero-weight vertices}
First, we arrange that every map in the support of the measure meets $Z$. Let $\rho$ be a measure supported on minimum maps of cost $L$. We use one of the following two sets of assumptions:
\begin{props}
\item\label{norm391:cutoff:1} all costs are nonnegative and all full maps cost more than $2L$;
\item\label{norm391:cutoff:2} for some $M$ with $1/r<M<1$, we have $L=1-M$, the vertices on maps in the support have costs $0\le c<M$, and all full maps cost at least one.
\end{props}
In the second case we sometimes also assume that every edge has total cost at least some $\kappa$ with
$\nu\kappa>1-2M$.

\begin{lemma}
\label{norm391:cutoff}\label{norm391:cutoff:rest}
Suppose that $\rho$ is supported on minimum maps of cost $L$ and satisfies~\eqref{norm391:counter}, and that~\ref{norm391:cutoff:1} or~\ref{norm391:cutoff:2} holds. Then there is a subdigraph $D'$ of $D$, a measure $\rho'$ of positive mass on maps into $D'$, and new costs $c'$ on $D'$, such that the following hold.
The minimum maps into $D'$ for $c'$ are exactly the minimum maps for $c$ which lie in $D'$, every map in the support of $\rho'$ is a minimum map which meets the set of vertices of cost zero for $c'$, and $\rho'$ satisfies~\eqref{norm391:counter} in $D'$ with the same $\varepsilon$. Moreover:
\begin{itemize}
\item in case~\ref{norm391:cutoff:1}, the new costs satisfy $0\le c'\le L'$, where $L'$ is the new minimum cost, and all full maps into $D'$ cost more than $2L'$;
\item in case~\ref{norm391:cutoff:2}, for some $M'$ with $M\le M'\le 1$, the new costs satisfy $0\le c'<M'$, the new minimum cost is $L'=1-M'$, and all full maps into $D'$ cost at least one. If every edge has cost at least $\kappa$, where $\nu\kappa>1-2M$, then every edge has new cost at least some $\kappa'$ with $\nu\kappa'>1-2M'$.
\end{itemize}
\end{lemma}
\begin{proof}
First, restrict $D$ to the set of vertices of positive multiplicity. Each of these vertices lies on a map of cost $L$, so it has cost at most $L$ in case~\ref{norm391:cutoff:1}, and less than $M$ in case~\ref{norm391:cutoff:2}. In particular, the completion sets of every minimum map are now disjoint, as a common vertex would complete it to a full map of cost at most $2L$ in case~\ref{norm391:cutoff:1}, and less than $L+M=1$ in case~\ref{norm391:cutoff:2}.

Let $\alpha$ be the maximum, over the maps $Q$ in the support of $\rho$, of the minimum cost of a vertex of $Q$. Keep only the maps in the support all of whose positions have cost at least $\alpha$, and let $D'$ be the subdigraph induced by the vertices of these maps. The kept maps have positive mass $m'$, and each of them has a position of cost exactly $\alpha$.
Now let $Q'$ be a minimum map into $D$ all of whose positions lie in $D'$. Neither completion set of $Q'$ contains a vertex of cost less than $\alpha$: inserting such a vertex and deleting the opposite endpoint, which has cost at least $\alpha$, would give a cheaper map. Each map in the support which was not kept has a position of cost less than $\alpha$, so, as the completion sets of $Q'$ are disjoint, it contributes at most $r-2$ positions to them. Subtracting these contributions from~\eqref{norm391:counter}, the multiplicity function $p'$ of the kept part of $\rho$ satisfies
\[
\Phi_{p'}(Q')\ge P-m+\varepsilon-(r-2)(m-m')=(r-2)m'+\varepsilon.
\]
Since the kept maps have cost $L$, the minimum maps into $D'$ are exactly the minimum maps into $D$ which lie in $D'$, and so~\eqref{norm391:counter} holds for $p'$ in $D'$.

In case~\ref{norm391:cutoff:1}, subtract $\alpha$ from every cost in $D'$. Each map has $r-1$ positions, so the new minimum cost is $L'=L-(r-1)\alpha\ge0$, the minimum maps are unchanged, and each vertex of $D'$, which lies on a kept map, has new cost at most $L'$. Full maps now cost more than $2L-r\alpha\ge2L'$. Finally, every kept map has a position of new cost zero.

In case~\ref{norm391:cutoff:2}, we have $\alpha\le L/(r-1)<1/r$, as $M>1/r$. Subtract $\alpha$ from every cost and divide by $1-r\alpha>0$, and put
\[
 M'=\frac{M-\alpha}{1-r\alpha}\qquad\text{and}\qquad
 L'=\frac{L-(r-1)\alpha}{1-r\alpha}=1-M'.
\]
As $rM>1$, we have $M'\ge M$, and $M'\le1$ as $L'\ge0$. All new costs are less than $M'$, and full maps still cost at least one. If edges have cost at least $\kappa$, then they have new cost at least $\kappa'=(\kappa-2\alpha)/(1-r\alpha)$, and
\[
 \nu\kappa'-(1-2M')
 =\frac{\nu\kappa-(1-2M)+2\alpha}{1-r\alpha}>0,
\]
using $r-2-2\nu=2$. Neither operation changes which maps are minimum, so the lemma follows.
\end{proof}

\subsection{Zero-weight positions in minimum maps}
We use the following classes of minimum maps, according to their endpoints. A minimum map is in the class $\mathcal E$ if both of its endpoints are zero endpoints. It is in the class $\mathcal A$ if it has a zero endpoint and an endpoint of positive cost, and the neighbour of the zero endpoint in the path $H-j$ is also a zero position, where $j$ is the deleted position. Finally, it is in the class $\mathcal C$ if it has a zero endpoint and an endpoint of positive cost, and the neighbour of the zero endpoint in $H-j$ has positive cost.
We also let
\[
 S=\{z\in Z:N^-(z)\cap Z=\varnothing\}\qquad\text{and}\qquad
 T=\{z\in Z:N^+(z)\cap Z=\varnothing\}
\]
be the vertices of $Z$ with no inneighbour in $Z$ and with no outneighbour in $Z$, respectively.

\begin{lemma}
\label{endpoint405:zero-endpoints}
Let $\rho$ be a measure on minimum maps with $p(Z)\ge2m$. Suppose that every minimum map has disjoint
completion sets, and that~\eqref{norm391:counter} holds. Then the following hold.
\begin{props}
\item If some minimum map lies in $\mathcal E\cup\mathcal A$, then $S$ or $T$ is avoided.\label{zeroend:1}
\item If every minimum map lies in $\mathcal E\cup\mathcal A$, then both
$S$ and $T$ are avoided, and so $p(S),p(T)\le m-\varepsilon$.\label{zeroend:2}
\item For each map in $\mathcal A$, the zero endpoint is internal. It lies in $S\setminus T$ if the edge between it and its neighbour in $H-j$ points away from it, and in $T\setminus S$ otherwise.\label{zeroend:3}
\end{props}
\end{lemma}
\begin{proof}
Recall that the completion pair of every minimum map has multiplicity on $Z$ at least $p(Z)-m+\varepsilon>m$, and that every avoided set has multiplicity at most $m-\varepsilon$.

We first prove~\ref{zeroend:3}. Let $Q\in\mathcal A$, with zero endpoint $e$, and let $e'$ be the neighbour of $e$ in $H-j$. By Observation~\ref{insertion:rule}, the completion set at $e$ is disjoint from $Z$. If $e$ were a turn, then the edges $ej$ and $ee'$ would point in the same direction relative to $e$, so the image of $e'$, which is in $Z$, would lie in the completion set at $e$, a contradiction. Thus $e$ is internal. If the edge between $e$ and $e'$ points from $e$ to $e'$, then the edge between $j$ and $e$ points into $e$, so the completion set at $e$ is the inneighbourhood of the image of $e$. Hence the image of $e$ has no inneighbour in $Z$, but has an outneighbour in $Z$, and so lies in $S\setminus T$. The other case is symmetric.

We also note how to turn a map in $\mathcal A$ into one in $\mathcal E$. The other completion set of $Q$ contains a vertex of $Z$, as the pair has positive multiplicity on $Z$. Inserting such a vertex and deleting the zero endpoint $e$ gives a minimum map whose endpoints are both zero positions, that is, a map in $\mathcal E$.

Now let $Q\in\mathcal E$, with deleted position $j$. If $j$ is a source of $H$, then the completion sets of $Q$ are the inneighbourhoods of two vertices of $Z$, which contain no vertex of $T$; thus $T$ is avoided. Similarly, if $j$ is a sink, then $S$ is avoided.
Suppose instead that $j$ is internal, with edges $j-1\to j\to j+1$ say, and let $u,v\in Z$ be the images of $j-1$ and $j+1$, so that the completion pair is $(N^+(u),N^-(v))$. Suppose that $u$ has an outneighbour $a\in Z$. Insert $a$ and delete $j+1$, obtaining a minimum map with completion pair $(N^+(a),W)$ say. No vertex of $N^+(a)$ lies in $S$, as $a\in Z$. Thus, if $W\cap S=\varnothing$, then $S$ is avoided. Otherwise, take $s\in W\cap S$, and replace $v$ by $s$ in $Q$. This gives a minimum map with completion pair $(N^+(u),N^-(s))$, which avoids $S$, as $u,s\in Z$ and $s$ has no inneighbour in $Z$.
Reversing all edges, we see similarly that $T$ is avoided if $v$ has an inneighbour in $Z$. As the completion pair of $Q$ has positive multiplicity on $Z$, $u$ has an outneighbour in $Z$ or $v$ has an inneighbour in $Z$, so $S$ or $T$ is avoided. Together with the previous paragraph, this proves~\ref{zeroend:1}.

For~\ref{zeroend:2}, suppose that every minimum map lies in $\mathcal E\cup\mathcal A$, but that $T$ is not avoided. By~\ref{zeroend:1}, $S$ is avoided, so $p(S)<m$.
Consider a map in $\mathcal A$ whose zero endpoint $x$ lies in $T$, so that, by~\ref{zeroend:3}, the edge between $x$ and its neighbour $y$ in $H-j$ is $y\to x$. For every vertex $a\in Z$ of the other completion set, inserting $a$ and deleting $x$ gives a map in $\mathcal E$ whose deleted position is internal, with endpoint images $y\to *\to a$. By the previous paragraph, if $a$ had an inneighbour in $Z$, then $T$ would be avoided. Hence $a\in S$, and so the other completion set has multiplicity at most $p(S)<m$ on $Z$, a contradiction. Thus, every map in $\mathcal A$ has its zero endpoint in $S$.

Now consider a map in $\mathcal E$, with deleted position $j$. By the arguments above, $j$ is not a source, and if $j$ is internal, with endpoint images $u\to*\to v$, then $v\in S$. If $j$ is a sink, insert an outneighbour $a\in Z$ of the image $u$ of one endpoint, and delete the other endpoint. If the new map lies in $\mathcal A$, then its zero endpoint $a$ has the inneighbour $u$ in $Z$, so it does not lie in $S$, which we have just excluded. Otherwise the new map lies in $\mathcal E$, and its deleted position is a source or is internal with second endpoint image $a$, which has the inneighbour $u\in Z$; either way, $T$ would be avoided.
Thus, every minimum map has a position mapped to $S$. As every map in the support of $\rho$ is minimum, this gives $p(S)\ge m$, a contradiction. Reversing all edges shows that $S$ is avoided as well, and this proves~\ref{zeroend:2}.
\end{proof}

When every vertex has cost zero, every map is minimum, and the lemma takes the following simpler form.

\begin{lemma}\label{base392:extremes}
Let $H$ have order $r\ge6$, and let $D_0$ be a digraph which receives no homomorphism from $H$.
Let $\rho$ be a measure of mass $m>0$ on maps into $D_0$, with multiplicity function
$p$, and let $P=(r-1)m$. Suppose that $\Phi_p(Q)\ge q$ for every map $Q$, including those outside the support of $\rho$, where $q>P-m$.
Then the set $S$ of vertices of $D_0$ with no inneighbour, and the set $T$ of vertices with no outneighbour, are each avoided by the completion pair of some map. In particular,
$p(S),p(T)\le P-q<m$.
\end{lemma}
\begin{proof}
Give every vertex cost zero. Then every map is minimum
and lies in $\mathcal E$, and $p(Z)=P\ge2m$.
The completion sets of every map are disjoint, as $D_0$ receives no homomorphism from $H$. Apply Lemma~\ref{endpoint405:zero-endpoints}
with $\varepsilon=q-P+m>0$.
\end{proof}

In the next three lemmas, we use either the \emph{strong cost assumptions}
\[
 0\le c\le L\quad\text{and all full maps cost more than }2L,
\]
or, when $r\ge6$, the \emph{weak cost assumptions}
\[
0<M\le1/2,\quad L=1-M,\quad 0\le c<M,\quad \nu\kappa>1-2M,
\]
all full maps cost at least one, and every edge costs at least $\kappa$.

\begin{lemma}
\label{norm391:boundary}\label{norm391:boundary:rest}
Let $\rho$ be a measure on minimum maps of cost $L$, each of which meets $Z$, and suppose that~\eqref{norm391:counter} holds and that the completion sets of every minimum map are disjoint. Then the following hold.
\begin{props}
\item\label{boundary:two} No minimum map has two endpoints of positive cost.
\item\label{boundary:one} Under the strong cost assumptions, or (when $r\ge6$) under the weak cost assumptions, no minimum map has exactly one zero position.
\item\label{boundary:C} Suppose that $r\ge6$ and $p(Z)\ge2m$. Then not every minimum map lies in $\mathcal C$.
\end{props}
\end{lemma}
\begin{proof}
For~\ref{boundary:two}, suppose that a minimum map has two endpoints of positive cost. Inserting a vertex of $Z$ from either completion set and deleting the opposite endpoint would give a cheaper map, so both completion sets are disjoint from $Z$. Their total multiplicity is then at most $P-p(Z)\le P-m$, as every map in the support of $\rho$ meets $Z$. This contradicts~\eqref{norm391:counter}.

For~\ref{boundary:one}, suppose that some minimum map has exactly one zero position. By~\ref{boundary:two}, this position is an endpoint. Removing it as well as the deleted position, we obtain a path of $r-2$ positions of positive cost, which we call the \emph{positive part} of the map. The two missing positions are consecutive, $t$ and $s$ say, where we choose the labels so that the edge between them is $t\to s$. Let $U$ and $V$ be the completion sets at the positions adjacent to $t$ and $s$ in the positive part. Thus, as far as the positive part is concerned, a vertex can be put at $t$ (respectively, $s$) exactly if it lies in $U$ (respectively, $V$).

We claim that there are minimum maps with this positive part which have a zero position at $t$ and the deleted position $s$, and also ones with a zero position at $s$ and the deleted position $t$. Indeed, our map has one of these forms. By Observation~\ref{insertion:rule}, its completion set at the zero endpoint is disjoint from $Z$, so, as its completion pair has positive multiplicity on $Z$, the other completion set contains a vertex of $Z$. Inserting this vertex and deleting the zero endpoint gives a map of the other form.

Among all such positive parts, choose one which maximises $\alpha=\min(A,B)$, where $A$ and $B$ are the costs of the positions of the positive part adjacent to $t$ and $s$, respectively. Fix minimum maps with this positive part and zero positions at $s$ and at $t$, with images $z_s$ and $z_t$. Their completion pairs are $(U,W_s)$ and $(V,W_t)$, where $W_s=N^-(z_s)$ and $W_t=N^+(z_t)$.
If $x\in W_s$, then putting $x$ at $t$ and deleting the position of cost $A$ next to $t$ gives a map, whose cost is $L-A+c(x)$. Therefore $c(x)\ge A$, and equality would give a minimum map with two positive endpoints, contradicting~\ref{boundary:two}. Thus $c>A$ on $W_s$, and similarly $c>B$ on $W_t$. Also, by Observation~\ref{insertion:rule}, $W_s$ and $W_t$ are disjoint from $Z$.
For each vertex $v$, let
\[
\delta(v)=2-\ind_U(v)-\ind_{W_s}(v)-\ind_V(v)-\ind_{W_t}(v),
\]
which is nonnegative as the two completion pairs are disjoint.

We next show that $U\cap V\cap Z=\varnothing$. Suppose first that $t$ and $s$ are both turns. In the map with $z_s$ at $s$, put the image of the position adjacent to $s$, which has cost $B$, also at $t$; as $s$ is a turn, this is compatible with the edge $t\to s$. Deleting the position of cost $A$ next to $t$ gives a map of cost $L-A+B$, so $B\ge A$, and equality would give two positive endpoints. Thus $B>A$, and, symmetrically, $A>B$, which is impossible. So one of $t$ and $s$ is internal, say $t$. If $z\in U\cap V\cap Z$, then put the image of the position next to $t$ at $t$, put $z$ at $s$, and delete the position next to $t$. As $t$ is internal, this is a map, and it is a minimum map with two positive endpoints, a contradiction. The case where $s$ is internal is symmetric. Therefore $U\cap V\cap Z=\varnothing$, and, as $W_s,W_t$ are disjoint from $Z$, we have $\delta\ge1$ on $Z$.

Now take any minimum map with exactly one zero position. Let $x$ be the cheaper of the two positions of its positive part next to its missing positions, and let $y$ be the neighbour of $x$ in its positive part. Then $c(x)\le\alpha$ by our choice of $\alpha$, so $x\notin W_s\cup W_t$.
Suppose that $\delta$ vanishes at the images of all positions of the positive part. Then $x\in U\cap V$. Suppose that $x\to y$; the other case is symmetric. As shown above, there is a version of this map in which $x$ is an endpoint, so, if $y$ had an inneighbour in $Z$, then replacing $x$ by this vertex would give a cheaper map. Hence $N^-(y)\cap Z=\varnothing$, and so $y\notin W_t$. As $\delta(y)=0$, we obtain $y\in V$. Now put $x$ at $t$ and $y$ at $s$ in the chosen positive part, giving a full map. Its cost is $L+c(x)+c(y)$.
Under the strong cost assumptions, this cost is more than $2L$, so $c(x)+c(y)>L$, which is impossible as $x$ and $y$ lie on a map of cost $L$. (Note that the positive part has at least two positions even when $r=4$.)
Under the weak cost assumptions, this cost is at least one, so $c(x)+c(y)\ge M$. The other $r-4$ positions of the positive part containing $x,y$ contain $\nu$ disjoint pairs of adjacent positions, each of cost at least $\kappa$. Thus $L\ge M+\nu\kappa>1-M=L$, again a contradiction.

Therefore, every minimum map with exactly one zero position has a position of positive cost with $\delta\ge1$ at its image, and so the images of its positions have total $\delta$ at least two. Minimum maps with at least two zero positions also have total $\delta$ at least two, as $\delta\ge1$ on $Z$. Every map in the support of $\rho$ is of one of these kinds, so integrating against $\rho$ gives $\sum_vp(v)\delta(v)\ge2m$. On the other hand, applying~\eqref{norm391:counter} to the two chosen minimum maps gives
\[
\sum_vp(v)\delta(v)=2P-\Phi_p(\text{first map})-\Phi_p(\text{second map})<2m,
\]
a contradiction.

For~\ref{boundary:C}, suppose that every minimum map lies in $\mathcal C$, and take one. Remove its zero endpoint and its deleted position, and call the two missing positions $j$ and $j+1$, so that the remaining positions are $j+2,\ldots,j-1$. Let $u$ and $v$ be the images of $j-1$ and $j+2$, which have positive cost, and let $U$ and $V$ be the corresponding completion sets for $j$ and $j+1$. As before, there are minimum maps with zero positions at $j$ and at $j+1$, and the completion set at the zero endpoint of each is disjoint from $Z$. As the two completion pairs have multiplicity at least $p(Z)-m+\varepsilon$ on $Z$, and are disjoint, this gives
\begin{equation}\label{boundary:UVZ}
 p(U\cap V\cap Z)\ge 2(p(Z)-m+\varepsilon)-p(Z)\ge 2\varepsilon>0.
\end{equation}

Suppose that $j$ is internal. Then we can move $u$ into $j$, put a vertex of $U\cap V\cap Z$ at $j+1$, and delete the old position $j-1$ of $u$. This gives a minimum map whose endpoint $j$ has positive cost, so, as it lies in $\mathcal C$, its other endpoint $j-2$ is a zero position. By Observation~\ref{insertion:rule}, the completion set of this map at $j$, which is a neighbourhood of $u$, has multiplicity at least $p(Z)-m+\varepsilon$ on $Z$. So does $U$, which is the neighbourhood of $u$ in the other direction, as $j$ is internal. Hence $p(N^+(u)\cap N^-(u)\cap Z)\ge2\varepsilon>0$, that is, $u$ forms a $2$-cycle with some vertex of $Z$. Symmetrically, if $j+1$ is internal, then $v$ forms a $2$-cycle with some vertex of $Z$.

We now show that we can move the two missing positions forward by two, keeping all positions outside $\{j,\ldots,j+3\}$ unchanged. If $j+1$ is internal, put a vertex of $U\cap V\cap Z$ at $j$, move $v$ to $j+1$, and delete $j+2$. This gives a minimum map whose endpoint $j+1$ has positive cost, so its other endpoint $j+3$ is forced to be a zero position.
If $j+1$ is a turn, we first show that $c(u)=c(v)$. Putting $v$ also at $j$, and a vertex of $V\cap Z$ at $j+1$, and deleting $j-1$, gives a map of cost $L-c(u)+c(v)$, so $c(v)\ge c(u)$. For the reverse inequality, if $j$ is a turn, we copy $u$ in the same way. If $j$ is internal, we put at $j$ a vertex of $Z$ forming a $2$-cycle with $u$, which exists by the previous paragraph, put $u$ at $j+1$, and delete $j+2$. The latter map has cost $L+c(u)-c(v)$.
Now put at $j$ a vertex of $U\cap Z$ if $j$ is a turn, or a vertex of $Z$ forming a $2$-cycle with $u$ if $j$ is internal, put $u$ at $j+1$, and delete $j+2$. This gives a minimum map of cost $L+c(u)-c(v)=L$ whose endpoint $j+1$ has positive cost. Again, its other endpoint $j+3$ is forced to be a zero position.
In every case, the new minimum map lies in $\mathcal C$, its zero endpoint is $j+3$, and the neighbour $j+4$ of $j+3$ has positive cost, by the definition of $\mathcal C$. So we can repeat this argument with $j+2$ in place of $j$.
After $r/2-1$ steps, the position $j-1$, which has positive cost and has not been changed, is forced to be a zero position, a contradiction. All the inequalities used concern the same function $p$, so this completes the proof.
\end{proof}

\begin{lemma}
\label{norm394:zero-extremes}\label{norm394:zero-extremes:rest}
Let $\rho$ be a measure on minimum maps of cost $L\ge0$ satisfying~\eqref{norm391:counter}. Suppose that the completion sets of every minimum map are disjoint, and that every minimum map has at least two zero positions and does not have two endpoints of positive cost.
Then every minimum map lies in $\mathcal E\cup\mathcal A$, the zero endpoint of each map in $\mathcal A$ is internal and lies in $S\setminus T$ or in $T\setminus S$ as described in Lemma~\ref{endpoint405:zero-endpoints}\ref{zeroend:3}, and both $S$ and $T$ are avoided, so that
\begin{equation}\label{norm391:extremal}
 p(S),p(T)\le m-\varepsilon.
\end{equation}
These assumptions hold if $L>0$, the strong cost assumptions hold, and every map in the support of $\rho$ meets $Z$.
\end{lemma}
\begin{proof}
We first check the last assertion. Under the strong cost assumptions, completion sets of minimum maps are disjoint, as a common vertex would complete a full map of cost at most $2L$.
Lemma~\ref{norm391:boundary} excludes minimum maps with exactly one zero position or with two positive endpoints, and a minimum map with no zero position has two positive endpoints, as $L>0$.

Now, under the assumptions of the lemma, every map in the support has at least two zero positions, so $z:=p(Z)\ge2m$.
A minimum map which is not in $\mathcal E\cup\mathcal A$ has a positive endpoint and a zero endpoint, and so lies in $\mathcal C$. If $r=4$, a map in $\mathcal C$ has only one zero position, so $\mathcal C$ is empty. Suppose that $r\ge6$ and that some minimum map lies in $\mathcal C$. As in the proof of Lemma~\ref{norm391:boundary}\ref{boundary:C}, let $t\to b$ be its two missing positions and let $U,V$ be the completion sets for $t$ and $b$, so that there are minimum maps with completion pairs $(U,N^-(z_b))$ and $(V,N^+(z_t))$, where $z_b,z_t\in Z$. By Observation~\ref{insertion:rule}, the completion sets $N^-(z_b)$ and $N^+(z_t)$ are disjoint from $Z$, and therefore
\[
 p(U\cap Z),\ p(V\cap Z)\ge z-m+\varepsilon>m.
\]
Every vertex $u\in U\cap Z$ can replace $z_t$, and then, again by Observation~\ref{insertion:rule}, $N^+(u)$ is disjoint from $Z$, that is, $u\in T$. Thus $U\cap Z\subseteq T$, and similarly $V\cap Z\subseteq S$. As avoided sets have multiplicity at most $m-\varepsilon$, neither $S$ nor $T$ is avoided. By Lemma~\ref{endpoint405:zero-endpoints}\ref{zeroend:1}, every minimum map then lies in $\mathcal C$, contradicting Lemma~\ref{norm391:boundary}\ref{boundary:C}.

Thus, every minimum map lies in $\mathcal E\cup\mathcal A$, and the remaining assertions follow from Lemma~\ref{endpoint405:zero-endpoints}\ref{zeroend:2} and~\ref{zeroend:3}.
\end{proof}

To deal with three or more zero positions, we use the structure of $D[Z]$ when it has no long walks. Suppose that $D[Z]$ has no directed walk with four edges. Let
$B=Z\setminus(S\cup T)$, and let $S_1$ and $T_1$ be the vertices of
$B$ with no inneighbour in $B$ and with no outneighbour in $B$, respectively.
Then $D[B]$ has no directed walk with two edges, since such a walk could be extended by one edge at each end within $Z$, as its endpoints lie in $B$.
Thus
\begin{equation}\label{norm391:four-cover}
 B=S_1\cup T_1,\quad Z=S\cup T\cup S_1\cup T_1,
 \quad N^-(S_1)\cap Z\subseteq S,
 \quad N^+(T_1)\cap Z\subseteq T.
\end{equation}

\begin{lemma}
\label{norm394:zero-strip}\label{norm394:zero-strip:rest}\label{norm398:threezero-exclusion}
Let $r\ge6$, and let $\rho$ be a measure on minimum maps of cost $L\ge0$ satisfying~\eqref{norm391:counter}. Suppose that the completion sets of every minimum map are disjoint, that $D[Z]$ has no directed walk with four edges, and that $S$ and $T$ are avoided. Suppose also that every minimum map lies in $\mathcal E\cup\mathcal A$, and that the zero endpoint of each map in $\mathcal A$ is internal and lies in $S\cup T$. Then the following hold.
\begin{props}
\item\label{zerostrip:three} If every minimum map has at least three zero positions, then no minimum map has exactly three.
\item\label{zerostrip:four} If $p(Z)\ge4m$, then $L=0$ and the sign sum of $H$ is zero.
\end{props}
\end{lemma}
\begin{proof}
Let $z=p(Z)$. In both parts we have $z\ge3m$, so the completion pair of every minimum map has multiplicity at least $z-m+\varepsilon\ge 2m+\varepsilon$ on $Z$, while every avoided set has multiplicity at most $m-\varepsilon$.
Recall that a map in $\mathcal A$ can be turned into a map in $\mathcal E$ with the same positions of positive cost, by inserting a vertex of $Z$ and deleting its zero endpoint.

We first note the following rule. Let $Q$ be a map in $\mathcal E$ whose deleted position is internal, with endpoint images $u\to*\to v$, and suppose that $v\in S_1$. Then $S_1$ is avoided. Indeed, $v\notin S$, so $v$ has an inneighbour $s\in Z$, and $s\in S$ by~\eqref{norm391:four-cover}. Inserting $s$ and deleting the endpoint mapped to $u$ gives a minimum map with completion pair $(W,N^-(s))$ say, and, as $N^-(s)\cap Z=\varnothing$, we have $p(W\cap Z)\ge z-m+\varepsilon>p(S)$. Choose $w\in(W\cap Z)\setminus S$, and replace $u$ by $w$ in $Q$. The resulting minimum map has completion pair $(N^+(w),N^-(v))$, which avoids $S_1$ by~\eqref{norm391:four-cover}. Reversing all edges, $T_1$ is avoided if $u\in T_1$.

For~\ref{zerostrip:three}, note first that if $L=0$, then all $r-1\ge5$ positions of every minimum map are zero positions. So we may assume that $L>0$.
We claim that, for a map in $\mathcal E$, the neighbours of the two endpoints in $H-j$ are not both of positive cost. Indeed, otherwise inserting a vertex of $Z$ and deleting an endpoint would give a map in $\mathcal A$, whose zero endpoint is the old deleted position, so the deleted position is internal, with endpoint images $u\to *\to v$ say. Inserting an outneighbour of $u$ in $Z$ and deleting the other endpoint gives a map in $\mathcal A$ whose zero endpoint has an inneighbour in $Z$, so this outneighbour lies in $T$. Similarly, every inneighbour of $v$ in $Z$ lies in $S$. Thus, the completion pair has multiplicity at most $p(S)+p(T)\le2m-2\varepsilon$ on $Z$, a contradiction.

Now consider a minimum map with exactly three zero positions; we may assume it lies in $\mathcal E$. Two of its zero positions are its endpoints, and, by the claim, the third is next to one of them. Thus its three zero positions and its deleted position form four consecutive positions $a,b,c,d$. (A map in $\mathcal A$ with exactly three zero positions is turned into a map in $\mathcal E$ with the same four positions.)
Below, for $\ell\in\{a,b,c,d\}$, a \emph{map missing $\ell$} is a minimum map with the same positions of positive cost as the given one and with deleted position $\ell$. Thus its other three positions in $\{a,b,c,d\}$ are zero positions.
As our map lies in $\mathcal E$, it misses $b$ or $c$.

Suppose that there is a map missing $b$, but none missing $c$. Then the completion set at $a$ of the map missing $b$ is disjoint from $Z$, as otherwise inserting a vertex of $Z$ and deleting $c$ would give a map missing $c$. Every vertex of $Z$ in the completion set at $c$ can be inserted, and deleting $a$ gives a map missing $a$, which lies in $\mathcal A$ with zero endpoint $b$. By assumption, the zero endpoints of these maps lie in $S\cup T$, so the completion pair of the map missing $b$ has multiplicity at most $p(S)+p(T)\le2m-2\varepsilon$ on $Z$, a contradiction. By symmetry, there are maps missing $b$ and maps missing $c$.

Suppose next that there is no map missing $a$. Then, as before, the completion set at $c$ of a map missing $b$ is disjoint from $Z$. In particular, $c$ is not a turn, as otherwise the image of $d$ would lie in this completion set. Reversing all edges if necessary, we may assume that $b\to c\to d$. Take a map missing $c$, with images $x_b$ and $x_d$ at $b$ and $d$. Every outneighbour of $x_b$ in $Z$ lies in $T$, since inserting it and deleting $d$ gives a map in $\mathcal A$ whose zero endpoint has an inneighbour in $Z$. Every inneighbour $y$ of $x_d$ in $Z$ lies in $S$, since otherwise we could put $y$ at $c$, an inneighbour of $y$ in $Z$ at $b$, and delete $a$, giving a map missing $a$. Thus, the completion pair of the map missing $c$ has multiplicity at most $2m-2\varepsilon$ on $Z$, a contradiction. By symmetry, there are maps missing $a$ and maps missing $d$.
Both of these lie in $\mathcal A$, with zero endpoints $b$ and $c$ respectively, so $b$ and $c$ are internal, and, after reversing all edges if necessary, we have $a\to b\to c\to d$.

Take a map missing $a$, with images $s\to x\to y$ at $b,c,d$. Its zero endpoint $b$ is mapped into $S\cup T$, and $s$ has the outneighbour $x\in Z$, so $s\in S$. Every inneighbour of $x$ in $Z$ can replace $s$, and must then also lie in $S$, so $x\in S_1$.
Let $U$ be the completion set of this map at the position before $a$. As $N^-(s)\cap Z=\varnothing$, we have $p(U\cap Z)\ge z-m+\varepsilon$.
Inserting any $u\in U\cap Z$ and deleting $b$ gives a map in $\mathcal E$ whose deleted position is internal, with endpoint images $u\to*\to x$. As $x\in S_1$, the rule above shows that $S_1$ is avoided.
No such $u$ lies in $T\cup T_1$: otherwise the completion pair $(N^+(u),N^-(x))$ of this map would have its multiplicity on $Z$ in $T\cup S$, by~\eqref{norm391:four-cover}, which is at most $2m-2\varepsilon$.
Hence $U\cap Z\subseteq S\cup S_1$ by~\eqref{norm391:four-cover}, and as both of these sets are avoided, $p(U\cap Z)\le2m-2\varepsilon$, a contradiction.
This proves~\ref{zerostrip:three}.

For~\ref{zerostrip:four}, suppose that $z\ge4m$. Suppose that some map in $\mathcal E$ has an internal deleted position with endpoint images $u\to*\to v$ and $v\in S_1$. Then, in the rule above, we have $p(W\cap Z)\ge3m+\varepsilon$, so we can choose $w\in W\cap Z$ outside the three avoided sets $S,T,S_1$, and then $w\in T_1$ by~\eqref{norm391:four-cover}. Replacing $u$ by $w$ gives a map in $\mathcal E$ to which the reversed rule applies, so $T_1$ is avoided too. But then $Z$ is covered by four avoided sets, so $z\le4(m-\varepsilon)$, a contradiction. Hence no such map exists, and, by reversal, no map in $\mathcal E$ has an internal deleted position with endpoint images $u\to*\to v$ and $u\in T_1$.

Next, if a minimum map has a zero endpoint mapped to $s\in S$ and the edge from this endpoint to its neighbour is $s\to v$ with $v\in S_1$, then this endpoint is a turn. Indeed, if it were internal, then its completion set $N^-(s)$ would be disjoint from $Z$, so the other completion set would contain a vertex $a$ of $Z$. Inserting $a$ and deleting the endpoint mapped to $s$ would then give a map in $\mathcal E$ with an internal deleted position and endpoint images $a\to*\to v$, which we have just excluded. The reversed statement holds for $u\in T_1$ and $t\in T$ with $u\to t$.

As $p(S)+p(T)\le2m-2\varepsilon$, the completion pair of every minimum map has positive multiplicity on $B$. Take a minimum map, insert a vertex of $B$ from its completion pair (from the completion set at the positive endpoint, for a map in $\mathcal A$), and delete the opposite zero endpoint. The new map cannot lie in $\mathcal A$, as its zero endpoint would lie in $B$, and not in $S\cup T$. So we obtain a map in $\mathcal E$ with an endpoint mapped into $B$.
Fix a directed path in $D[Z]$ through the image of this endpoint, of length two or three, as follows. If the image is isolated in $D[B]$, take an inneighbour in $S$ and an outneighbour in $T$. Otherwise it is joined to a vertex of $B$, and, as $D[B]$ has no walk with two edges, we obtain a path whose two inner vertices lie in $S_1$ and $T_1$, and whose outer vertices lie in $S$ and $T$.

Now move the deleted position forward by one. That is, fill it with the predecessor or successor, on this fixed path, of the image of the adjacent endpoint, according to the direction of the edge of $H$ between them, and delete the other endpoint. The new endpoint which was not changed must be a zero position, as otherwise the new map would lie in $\mathcal A$. This is impossible both for a zero endpoint mapped to an inner vertex of the path, which lies in $B$, and for one mapped to an outer vertex, by the turn rule above. So the new map lies in $\mathcal E$. When the new vertex is an outer vertex of the path, the corresponding position of $H$ is a turn, so the next step goes back into the path, and the operation never leaves the path.
If $L>0$, continue until the first position which originally had positive cost. This position is unchanged, but is forced to be a zero position, a contradiction.
If $L=0$, continue indefinitely. After every $r$ steps, the position on the path has changed by the sign sum of $H$. As the path is bounded, the sign sum is zero.
\end{proof}

\subsection{Endpoint bounds}
We can now prove the endpoint bound in the three cost settings needed below.
Let $H$ have even order $r\ge6$, and let $\rho$ be a measure of mass $m>0$ on maps into $D$, with multiplicity function $p$.
We consider the following three sets of assumptions.
\begin{props}
\item\label{norm391:positive-endpoint:1} The costs satisfy $0\le c\le L$ for some $L>0$, all maps cost
at least $L$, and all full maps cost more than $2L$. The measure $\rho$ is supported
on maps of cost $L$ which meet $Z$, and $D[Z]$ has no directed walk with four edges. (The last condition holds if, for some position $j$, the span of $H-j$ is at most four, as then a walk with four edges in $D[Z]$ would receive a map of cost zero.)

\item\label{norm391:positive-endpoint:2} For some $0<M\le1/2$, put $L=1-M$. The costs satisfy $0\le c<M$,
all maps cost at least $L$, and all full maps cost at least one.
Every edge costs at least $\kappa$, where $\nu\kappa>1-2M$.
The measure $\rho$ is supported on maps of cost $L$ which meet $Z$.

\item\label{norm391:positive-endpoint:3} $H$ satisfies~\eqref{norm391:range}, every vertex has cost
zero, and $D$ receives no homomorphism from $H$. The measure $\rho$
may be supported on arbitrary maps.
\end{props}

\begin{lemma}
\label{norm391:positive-endpoint}\label{norm391:weak-endpoint}\label{norm391:zero-endpoint}
Under any of~\ref{norm391:positive-endpoint:1}--\ref{norm391:positive-endpoint:3}, some minimum map $Q$ satisfies $\Phi_p(Q)\le(r-2)m$.
\end{lemma}
\begin{proof}
Suppose otherwise, so that~\eqref{norm391:counter} holds for some $\varepsilon>0$.

\emph{Case~\ref{norm391:positive-endpoint:1}.}
By Lemma~\ref{norm394:zero-extremes}, every minimum map lies in $\mathcal E\cup\mathcal A$, and~\eqref{norm391:extremal} holds. Let $z=p(Z)$, so that $z\ge2m$.
We first show that no minimum map has exactly two zero positions. Suppose that one does. Turning it into a map in $\mathcal E$ if necessary, we may assume that both its endpoints are its zero positions. Inserting a vertex of $Z$ and deleting an endpoint gives a map in $\mathcal A$ whose zero endpoint is the old deleted position, so this position is internal. Thus we can label the two missing positions and the deleted position as $a\to b\to c$, where $b$ is deleted, and we let $A$ and $C$ be the costs of the positions adjacent to $a$ and $c$ outside $\{a,b,c\}$.
Let $u,v\in Z$ be the images of $a$ and $c$. Every outneighbour of $u$ in $Z$ lies in $T$, and every inneighbour of $v$ in $Z$ lies in $S$, as inserting it and deleting the other endpoint gives a map in $\mathcal A$. The completion pair $(N^+(u),N^-(v))$ has multiplicity at least $z-m+\varepsilon>m$ on $Z$, so, by~\eqref{norm391:extremal}, $u$ has an outneighbour in $Z$ and $v$ has an inneighbour in $Z$. Inserting these, we obtain minimum maps with the same positions of positive cost and deleted positions $c$ and $a$, respectively.
Let $U$ and $V$ be the completion sets for $a$ and $c$ defined by their neighbours outside $\{a,b,c\}$. By Observation~\ref{insertion:rule}, the completion sets at the zero endpoints of these two maps are disjoint from $Z$, so $p(U\cap Z),p(V\cap Z)\ge z-m+\varepsilon$, and
$p(U\cap V\cap Z)\ge z-2m+2\varepsilon>0$.
For $x\in U\cap V\cap Z$, putting $x$ at both $a$ and $c$ gives a map in $\mathcal E$ with deleted position $b$, and, by the same argument, $x$ has an inneighbour $s\in S$ and an outneighbour $t\in T$.
Thus, there are three minimum maps with the same positions of positive cost, with completion pairs
\[
 (U,N^-(s)),\quad(N^+(x),N^-(x)),\quad(V,N^+(t)).
\]
Let $\delta$ be three minus the sum of the indicator functions of these six sets, which is nonnegative as the three pairs are disjoint.
On $Z$, the sets $N^-(s)$ and $N^+(t)$ are empty, $N^+(x)\cap Z\subseteq T$ and $N^-(x)\cap Z\subseteq S$, and every vertex of $U\cap V\cap Z$ lies in $B$ by the argument just given. Hence $\delta\ge1$ on $Z$.
As in the proof of Lemma~\ref{norm391:boundary}, we have $c>A$ on $N^-(s)$ and $c>C$ on $N^+(t)$, since equality would give two positive endpoints.

Choose the positions of positive cost, among all minimum maps with exactly two zero positions, to maximise $\min(A,C)$. Consider any minimum map with exactly two zero positions, and suppose that $\delta$ vanishes at the images of all its positions of positive cost. Let $y$ be the cheaper of its two positions of positive cost next to its missing positions. Then $y\notin N^-(s)\cup N^+(t)$, so $\delta(y)=0$ gives $y\in U\cap V$. Let $z'$ be the neighbour of $y$ among its positions of positive cost, and suppose that $y\to z'$. As before, there is a version of this map in which $y$ is an endpoint, so $x\to z'$ is impossible, as replacing $y$ by $x$ would give a cheaper map. Thus $z'\notin N^+(x)$, and $\delta(z')=0$ gives $z'\in N^-(x)$. Putting $y,z',x$ at $a,b,c$ in the chosen map gives a full map of cost $L+c(y)+c(z')\le2L$, a contradiction. If $z'\to y$, we use $x,z',y$ instead.
Therefore, every minimum map with exactly two zero positions has total $\delta$ at least three over the images of its positions, and so does every minimum map with at least three zero positions. Integrating against $\rho$ gives $\sum_vp(v)\delta(v)\ge3m$, which contradicts the sum of the three inequalities~\eqref{norm391:counter} for the three chosen maps.
Thus every minimum map has at least three zero positions. By Lemma~\ref{norm394:zero-strip}\ref{zerostrip:three}, none has exactly three, so every minimum map has at least four, and $z\ge4m$. As $L>0$, this contradicts Lemma~\ref{norm394:zero-strip}\ref{zerostrip:four}.

\emph{Case~\ref{norm391:positive-endpoint:2}.}
Completion sets of minimum maps are disjoint,
as $L+c<1$. As $\nu\kappa>1-2M\ge0$, we have $\kappa>0$, so $Z$
is independent. By Lemma~\ref{norm391:boundary}, no minimum map has two positive endpoints or exactly one zero position, so $p(Z)\ge2m$. If a minimum map had two zero endpoints, then its completion sets would be neighbourhoods of vertices of $Z$, and so disjoint from $Z$, and their multiplicity would be at most $P-p(Z)<P-m$. Thus every minimum map has one positive endpoint and one zero endpoint, whose neighbour has positive cost as $Z$ is independent, that is, every minimum map lies in $\mathcal C$. This contradicts Lemma~\ref{norm391:boundary}\ref{boundary:C}.

\emph{Case~\ref{norm391:positive-endpoint:3}.}
Restrict $D$ to the vertices of positive multiplicity. Every map is minimum, its completion sets are disjoint, and $Z=V(D)$.
Let $S$ and $T$ be the sets of vertices with no inneighbour and with no outneighbour, respectively, and let $B=V(D)\setminus(S\cup T)$.
By Lemma~\ref{base392:extremes}, with $q=P-m+\varepsilon$, we have
$p(S),p(T)\le m-\varepsilon$, and both sets are avoided.

The digraph $D$ has a directed walk with four edges. Otherwise, Lemma~\ref{norm394:zero-strip}\ref{zerostrip:four} applies with $L=0$, $Z=V(D)$ and
$p(Z)=P=(r-1)m\ge4m$: all maps lie in $\mathcal E$, and $S,T$ are avoided. As the sign sum of $H$ is two, this is a contradiction.

As every orientation of an even cycle maps to a $2$-cycle, $D$ is an oriented graph, and, by Lemma~\ref{norm391:height}\ref{norm391:height:3}, its underlying graph is
$K_{\Delta+3}$-free. Let $\ell=\Delta+1$, and let $A$ be the set of
vertices which are internal vertices of some directed walk with $\ell$ edges. For every
$a\in A$, the map given by Lemma~\ref{norm391:height}\ref{norm391:height:2} has completion sets $N^-(a)$ and $N^+(a)$, so
\[
 d_p(a):=p(N^-(a))+p(N^+(a))\ge q.
\]
Choose a subwalk $z_0\to\cdots\to z_\ell$ of a walk with four edges, and let $i=\Delta-1\in\{1,2\}$, so that $z_i\in A$.
We claim that $z_i$ has no neighbour in $B\setminus A$. If $b\to z_i$ with $b\in B$, then $b$ has an inneighbour, and following it by $b$ and then by $z_i,\ldots,z_\ell$ gives a walk whose first $\ell$ edges have $b$ as an internal vertex. If $z_i\to b$ with $b\in B$, then $b$ has an outneighbour, and appending $b$ and this outneighbour to $z_0\cdots z_i$ gives a walk with $i+2=\ell$ edges in which $b$ is internal. (Repeated vertices do not affect either argument.)
As $D$ is oriented, the bound $d_p(z_i)\ge q$ gives
$p(B\setminus A)\le P-q=m-\varepsilon$, and hence
$p(A)\ge P-3(m-\varepsilon)$.

By Lemma~\ref{general:weighted-turan}, applied to $D[A]$, some $a\in A$ with $p(a)>0$ has at most $(1-1/(\Delta+2))p(A)$ multiplicity in its neighbourhood inside $A$, and at most $P-p(A)$ outside $A$. Thus
\[
 q\le d_p(a)\le P-\frac{p(A)}{\Delta+2},\qquad\text{so}\qquad
 p(A)\le(\Delta+2)(m-\varepsilon).
\]
Combining these bounds gives $P\le(\Delta+5)(m-\varepsilon)
<(\Delta+5)m\le(r-1)m=P$, by~\eqref{norm391:range}, a contradiction.
\end{proof}

Combining Lemma~\ref{norm391:positive-endpoint} with the restriction in Lemma~\ref{norm391:cutoff}, we obtain the following endpoint bound.

\begin{theorem}
\label{norm391:endpoint}
Let $H$ satisfy~\eqref{norm391:range}, and let $\rho$ be a measure of mass $m>0$ on
minimum maps of cost $L$. Suppose that one of the following holds:
\begin{props}
\item\label{norm391:endpoint:1} the costs are nonnegative, and every full map costs more than $2L$;
\item\label{norm391:endpoint:2} for some $M$ with $1/r<M<1$, we have $L=1-M$, the vertices on maps in the support have costs $0\le c<M$,
every full map costs at least one, and every edge costs at least $\kappa$, where
$\nu\kappa>1-2M$.
\end{props}
Then some minimum map $Q$ has $\Phi_p(Q)\le(r-2)m$.
\end{theorem}
\begin{proof}
Suppose otherwise, so that~\eqref{norm391:counter} holds for some $\varepsilon>0$, and apply Lemma~\ref{norm391:cutoff}. As the minimum maps after the restriction are minimum maps before it, we obtain a measure for which~\eqref{norm391:counter} holds, and we derive a contradiction by applying Lemma~\ref{norm391:positive-endpoint} to it, as follows.

In case~\ref{norm391:endpoint:1}, if the new minimum cost is zero, then all vertices have cost zero and the digraph receives no full map, which contradicts Lemma~\ref{norm391:positive-endpoint} in case~\ref{norm391:positive-endpoint:3}.
If the new minimum cost is positive, we use case~\ref{norm391:positive-endpoint:1}, as, by Lemma~\ref{norm391:height}\ref{norm391:height:1}, the span of $H-j$ is at most $\Delta+1\le4$ for some position $j$.
In case~\ref{norm391:endpoint:2}, if $M'\le1/2$, we use case~\ref{norm391:positive-endpoint:2}.
If $M'>1/2$, then full maps cost at least one, which is more than $2(1-M')$, so we are in case~\ref{norm391:endpoint:1} (this includes $M'=1$), which we have already treated.
\end{proof}

We can now deduce the bounds on optimal covers which we use for the remaining orientations.

\begin{theorem}
\label{norm391:normalized-cover}
Let $H$ satisfy~\eqref{norm391:range}, and let $F$ be a digraph of order $N$ with
$d:=\delta^0(F)>N/2$. Let $w$ be a deficient fractional
$H$-cover of $F$ of minimum total weight, and let
$M=\max w$ and $\kappa=\min_{xy\in E(F)}(w(x)+w(y))$.
Then
\[
 \min w=0,\qquad \frac2r\le M\le\frac12,\qquad
 \left(\frac r2-2\right)\kappa\le1-2M.
\]
Moreover, every partial homomorphism into $F$ costs at least $1-M$, and every vertex forming a $2$-cycle with a vertex of weight zero has weight at least $2/r$.
\end{theorem}
\begin{proof}
By Lemma~\ref{general:LP}\ref{LP:normal}, we have $0\le w\le1$ and $\min w=0$.
The two completion sets of any partial homomorphism have at
least $d$ vertices each, so they intersect, and a common vertex, of weight at most
$M$, completes it to a homomorphism of cost at least one. Hence every partial homomorphism costs at least $1-M$.
If $w(z)=0$, then $z$ forms a $2$-cycle with some vertex $t$, as $2d>N$, and mapping the two bipartition classes of $H$ to $z$ and $t$ gives $(r/2)w(t)\ge1$. In particular, $M\ge2/r>1/r$.

Suppose, for contradiction, that $\nu\kappa>1-2M$. Let $B_M$ be the set of vertices of weight $M$.
We claim that a tight homomorphism cannot have two distinct positions mapped to $B_M$. Indeed, removing these two positions from $H$ leaves at most two paths with $r-2$ vertices in total, which contain at least $(r-4)/2=\nu$ disjoint pairs of adjacent positions, each of cost at least $\kappa$. So such a homomorphism would cost at least $2M+\nu\kappa>1$.
Therefore, Lemma~\ref{general:max-deletion} applies, and gives a measure of mass $|B_M|$ on minimum maps of cost $L=1-M$ in $D=F-B_M$, for which every minimum map $Q$ in $D$ has $\Phi_p(Q)>(r-2)|B_M|$.
In $D$, all weights are less than $M$, and full maps still cost at least one.
If $M>1/2$, then full maps cost more than $2L$, and Theorem~\ref{norm391:endpoint}\ref{norm391:endpoint:1} gives a contradiction. Otherwise, Theorem~\ref{norm391:endpoint}\ref{norm391:endpoint:2} applies,
with our assumption $\nu\kappa>1-2M$, and again gives a contradiction.

Therefore $\nu\kappa\le1-2M$. As $\kappa\ge0$, this also gives $M\le1/2$.
\end{proof}
